% Auto-generated from book/09_quaternion_algebras.md — do not edit by hand
% Regenerated by scripts/md_to_latex.py

\chapter{Chapter 9 --- Quaternion Algebras and Ideal Theory}
\label{ch:ch09_algebras}

This chapter supplies the rigorous algebraic foundation for the class-group analogues and composition laws of Chapter 8. We introduce quaternion algebras over \(\mathbb{Q}\), their orders (especially the Hurwitz order), and the ideal theory that turns equivalence classes of ideals into a group. This framework is the natural home for making the Model constructions of previous chapters more arithmetic and for attacking Open Problem 6 with algebraic tools.

\textbf{Learning goals}

\begin{enumerate}
\item Define quaternion algebras over \(\mathbb{Q}\) and understand ramification.  
\item Work with orders (Lipschitz, Hurwitz, maximal) and their ideal theory.  
\item See how ideal classes form a group and how this relates to the class-group analogues of Chapter 8.  
\item Connect ideal theory back to composition of forms / flywheels and to Magic Islands (\textbf{Model} dictionary).  
\item Prepare algebraic language for rigorous uniqueness and composition statements (OP6 and beyond).  
\item Separate \textbf{algebraic}, \textbf{sequential-label}, and \textbf{angle-lock} invariants when discrete moduli label continuous geometric orbits (supporting stack \texttt{vortex\_\allowbreak{}math}).
\end{enumerate}

\textbf{Figures in this chapter}

\begin{center}
\small
\setlength{\tabcolsep}{4pt}
\begin{tabularx}{\textwidth}{@{}>{\raggedright\arraybackslash}p{0.12\textwidth}>{\raggedright\arraybackslash}X>{\raggedright\arraybackslash}X@{}}
\toprule
Tag & File & Role \\
\midrule
Fig. 9.1 & \texttt{figures/\allowbreak{}fig9\_\allowbreak{}1\_\allowbreak{}quaternion\_\allowbreak{}algebra.\allowbreak{}png} & Quaternion algebra ramification diagram \\
Fig. 9.2 & \texttt{figures/\allowbreak{}fig9\_\allowbreak{}2\_\allowbreak{}hurwitz\_\allowbreak{}order.\allowbreak{}png} & Hurwitz order units / comparison \\
Fig. 9.3 & \texttt{figures/\allowbreak{}fig9\_\allowbreak{}3\_\allowbreak{}ideal\_\allowbreak{}class\_\allowbreak{}group.\allowbreak{}png} & Ideal class group schematic \\
Fig. 9.4 & \texttt{figures/\allowbreak{}fig9\_\allowbreak{}4\_\allowbreak{}ideal\_\allowbreak{}to\_\allowbreak{}flywheel.\allowbreak{}png} & Ideal class \(\rightarrow\) flywheel (Model bridge) \\
Fig. 9.5 & \texttt{figures/\allowbreak{}fig9\_\allowbreak{}5\_\allowbreak{}modulus\_\allowbreak{}invariants.\allowbreak{}png} & Three-layer modulus invariants (algebra \(\cdot\) flow \(\cdot\) lock control) \\
Aux A9.1 & \texttt{figures/\allowbreak{}aux9\_\allowbreak{}1\_\allowbreak{}ramification\_\allowbreak{}table.\allowbreak{}png} & Hilbert-symbol ramification table \\
\bottomrule
\end{tabularx}
\end{center}


\textbf{Claim discipline}

\begin{center}
\small
\setlength{\tabcolsep}{4pt}
\begin{tabularx}{\textwidth}{@{}>{\raggedright\arraybackslash}X>{\raggedright\arraybackslash}X@{}}
\toprule
Claim & Type \\
\midrule
Quaternion algebras, Hilbert symbols, Hurwitz Euclidean property, Hurwitz class number 1 & \textbf{Theorem} (classical) \\
Dictionary between ideal classes and flux topographs / Magic Islands & \textbf{Model} \\
Separation of modulus invariants on irrational rotations (algebra vs progression vs angle lock) & \textbf{Model} (experimentally controlled; not a theorem of TN) \\
\texttt{qga/\allowbreak{}lib/\allowbreak{}quaternion\_\allowbreak{}algebra.\allowbreak{}py} helpers; external \texttt{vortex\_\allowbreak{}math} metrics & \textbf{Software fact} \\
\bottomrule
\end{tabularx}
\end{center}


\bigskip
\noindent\rule{\textwidth}{0.4pt}
\bigskip

\section{9.1 Quaternion algebras over \(\mathbb{Q}\)}
\label{ch:ch09_algebras:9-1-quaternion-algebras-over-mathbb-q}

A \textbf{quaternion algebra} over \(\mathbb{Q}\) is a 4-dimensional central simple algebra over \(\mathbb{Q}\). Up to isomorphism, each is given by square-free integers \(a,b\neq 0\) via

\[
\Bigl(\frac{a,b}{\mathbb{Q}}\Bigr)
=
\mathbb{Q}+\mathbb{Q}i+\mathbb{Q}j+\mathbb{Q}ij,
\qquad
i^2=a,\; j^2=b,\; ij=-ji.
\]

The Hamilton algebra corresponds (over \(\mathbb{R}\)) to \(a=b=-1\); as an algebra over \(\mathbb{Q}\) one studies \(\bigl(\frac{-1,-1}{\mathbb{Q}}\bigr)\).

\subsection{Ramification}
\label{ch:ch09_algebras:ramification}

At each place \(p\) (finite prime or \(\infty\)), the local Hilbert symbol \((a,b)_p\in\{\pm 1\}\) detects whether the algebra is a division algebra (\(-1\), \textbf{ramified}) or isomorphic to \(M_2\) (\textbf{split}). The set of ramified places is finite and of even cardinality; it classifies the algebra.

\begin{lstlisting}[style=qga]
QuaternionAlgebra(a, b)
  .ramified_places()
  .is_definite()          # True iff ramified at ?
  .hilbert_at(p)
\end{lstlisting}

\begin{figure}[htbp]
  \centering
  \includegraphics[width=0.92\textwidth,height=0.42\textheight,keepaspectratio]{fig9_1_quaternion_algebra.png}
  \caption{Figure 9.1 --- Ramification of a quaternion algebra.}
  \label{fig:fig9-1-quaternion-algebra}
\end{figure}

\begin{quote}\small\textit{Figure 9.1.* Example \(\bigl(\frac{-1,-1}{\mathbb{Q}}\bigr)\): ramified at \(2\) and \(\infty\) (definite over \(\mathbb{R}\)).}\end{quote}

\begin{figure}[htbp]
  \centering
  \includegraphics[width=0.92\textwidth,height=0.42\textheight,keepaspectratio]{aux9_1_ramification_table.png}
  \caption{Auxiliary Figure A9.1 --- Hilbert-symbol table.}
  \label{fig:aux9-1-ramification-table}
\end{figure}

\begin{quote}\small\textit{Auxiliary Figure A9.1.* Hilbert symbols for several \((a,b)\) at small places; red \(-1\) marks ramification.}\end{quote}

\textbf{Claim type.} Classification by ramification sets and Hilbert symbols: \textbf{Theorem} (classical). Implementation in the book helper: \textbf{Software fact} (pedagogical, not a full number-theory library).

\bigskip
\noindent\rule{\textwidth}{0.4pt}
\bigskip

\section{9.2 Orders in quaternion algebras}
\label{ch:ch09_algebras:9-2-orders-in-quaternion-algebras}

An \textbf{order} is a subring that is a full-rank \(\mathbb{Z}\)-lattice in the algebra. In the Hamilton setting:

\begin{center}
\small
\setlength{\tabcolsep}{4pt}
\begin{tabularx}{\textwidth}{@{}>{\raggedright\arraybackslash}X>{\raggedright\arraybackslash}X>{\raggedright\arraybackslash}X>{\raggedright\arraybackslash}X@{}}
\toprule
Order & Definition & Units & Euclidean? \\
\midrule
Lipschitz & \(\mathbb{Z}[i,j,k]\) & 8 & No \\
Hurwitz & all-integer or all-half-integer coords & 24 & Yes (maximal) \\
\bottomrule
\end{tabularx}
\end{center}


The Hurwitz order is preferred for arithmetic (Chapter 1 \textbf{Model} preference for lattice work; Euclidean algorithm is classical \textbf{Theorem}). Its unit group is the binary tetrahedral group---the set \(\Lambda_0\) of Chapter 3.

\begin{figure}[htbp]
  \centering
  \includegraphics[width=0.92\textwidth,height=0.42\textheight,keepaspectratio]{fig9_2_hurwitz_order.png}
  \caption{Figure 9.2 --- Hurwitz order.}
  \label{fig:fig9-2-hurwitz-order}
\end{figure}

\begin{quote}\small\textit{Figure 9.2.* 24 Hurwitz units in stereographic \(S^3\), and Lipschitz vs Hurwitz comparison.}\end{quote}

\begin{lstlisting}[style=qga]
HurwitzOrder().units          # 24
HurwitzOrder().is_euclidean() # True (classical fact)
HurwitzOrder().is_maximal()   # True (classical fact)
LipschitzOrder().n_units()    # 8
\end{lstlisting}

In general quaternion algebras, \textbf{maximal orders} play the role of rings of integers in number fields; different maximal orders need not be conjugate when the class number is nontrivial.

\bigskip
\noindent\rule{\textwidth}{0.4pt}
\bigskip

\section{9.3 Ideal theory and class groups}
\label{ch:ch09_algebras:9-3-ideal-theory-and-class-groups}

One studies left ideals, right ideals, and two-sided ideals of an order. For \textbf{definite} quaternion algebras (ramified at \(\infty\)), the left ideal class set of a maximal order is finite. Ideal multiplication induces a group structure on (two-sided) ideal classes, and related class sets control representation of integers by norms of ideals.

\subsection{Hurwitz class number (classical)}
\label{ch:ch09_algebras:hurwitz-class-number-classical}

The left ideal class number of the Hurwitz order is \textbf{1}: every left ideal is principal. The book helper reports this as a cited classical fact with a finite sample check of small-norm elements---not a full ideal enumeration.

\begin{lstlisting}[style=qga]
left_ideal_class_group(HurwitzOrder())
  -> IdealClassGroupResult(order=1, method='classical_hurwitz_class_number_one', ...)
\end{lstlisting}

\begin{figure}[htbp]
  \centering
  \includegraphics[width=0.92\textwidth,height=0.42\textheight,keepaspectratio]{fig9_3_ideal_class_group.png}
  \caption{Figure 9.3 --- Ideal class group schematic.}
  \label{fig:fig9-3-ideal-class-group}
\end{figure}

\begin{quote}\small\textit{Figure 9.3.* Ideal classes with multiplication. For Hurwitz, a single class; general definite orders may have larger finite class numbers.}\end{quote}

\textbf{Claim type.} Finiteness for definite maximal orders; Hurwitz class number 1: \textbf{Theorem} (classical). Full algorithmic enumeration of ideals for general \((a,b)\): out of scope for this sandbox.

\bigskip
\noindent\rule{\textwidth}{0.4pt}
\bigskip

\section{9.4 From ideals to flux topographs and flywheels (Model bridge)}
\label{ch:ch09_algebras:9-4-from-ideals-to-flux-topographs-and-flywheels-model-bridg}

Chapter 8 approximated class groups with composition of flux topographs. The natural dictionary is still a \textbf{Model}:

\begin{center}
\small
\setlength{\tabcolsep}{4pt}
\begin{tabularx}{\textwidth}{@{}>{\raggedright\arraybackslash}X>{\raggedright\arraybackslash}X@{}}
\toprule
Algebraic object & Geometric / Model object \\
\midrule
Ideal in a quaternion order & Flux configuration / supporting cycle on the gauged Hopf lattice \\
Ideal class & Equivalence class of reduced flux topographs \\
Ideal class group & Class-group analogue of Chapter 8 \\
Norm of an ideal & Stability score or \texttt{magic\_\allowbreak{}island\_\allowbreak{}score} \\
Ideal multiplication & Composition of flywheels (when well-defined --- OP6) \\
\bottomrule
\end{tabularx}
\end{center}


\begin{figure}[htbp]
  \centering
  \includegraphics[width=0.92\textwidth,height=0.42\textheight,keepaspectratio]{fig9_4_ideal_to_flywheel.png}
  \caption{Figure 9.4 --- From ideal class to flux flywheel.}
  \label{fig:fig9-4-ideal-to-flywheel}
\end{figure}

\begin{quote}\small\textit{Figure 9.4.* Schematic Model bridge. Not a proved functor.}\end{quote}

\textbf{Open Problem 6 (continued).} Use ideal theory of quaternion algebras to define a rigorous composition law on classes of flux configurations that is associative, gauge-compatible, and reduces to Gauss composition in suitable limits. Determine whether the resulting class group predicts Magic Island location and size.

\begin{lstlisting}[style=qga]
form_ideal_dictionary_entry()   # static dictionary for labs
class_group_analogue(...)       # Ch. 8 Model side
left_ideal_class_group(...)     # algebraic side (Hurwitz: order 1)
\end{lstlisting}

\bigskip
\noindent\rule{\textwidth}{0.4pt}
\bigskip

\section{9.5 Modulus, invariants, and angle--label locking}
\label{ch:ch09_algebras:9-5-modulus-invariants-and-angle-label-locking}

*Supporting stack: \href{https://github.com/kinaar8340/vortex_math}{\texttt{vortex\_\allowbreak{}math}} \(\cdot\) research note: \href{../notes/RESEARCH_NOTE_moduli.md}{\texttt{notes/\allowbreak{}RESEARCH\_\allowbreak{}NOTE\_\allowbreak{}moduli.\allowbreak{}md}}.*

Chapters 3--8 build discrete geometric objects (gauged Hopf lattice, flux topographs) whose \textbf{labels} and \textbf{equivalence classes} carry arithmetic meaning. Before those dictionaries become too ambitious, it is worth isolating a cleaner, lower-dimensional lesson: when symbolic patterns of Vortex Math---digital roots, the doubling circuit, and the privileged status of 3-6-9---are placed onto points generated by \emph{continuous geometry} rather than discrete integer sequences alone, the choice of labeling modulus does \textbf{not} act as a single uniform operator. It modifies \emph{different invariants} depending on how it is attached to the system.

\subsection{Geometric stage}
\label{ch:ch09_algebras:geometric-stage}

Consider an irrational rotation on the unit circle whose step size is a fixed multiple of the circle constant:

\[
\theta_k = k \cdot \frac{9}{\pi} \pmod{2\pi}.
\]

Because \(9/\pi\) is an irrational multiple of \(2\pi\), the orbit is dense: successive points come arbitrarily close to every location on the circle without closing into a finite regular polygon. The factor \textbf{9} echoes the digital-root base of Vortex Math; \textbf{?} is the circle's own constant. Geometry and numerology share the stage, but they are no longer forced to play the same role. (A coupled step \(\Delta\theta = m/\pi\) will appear below when the modulus is allowed to set the winding rate.)

\subsection{Three layers of invariant}
\label{ch:ch09_algebras:three-layers-of-invariant}

\textbf{Algebraic structure of the doubling map.} Changing the modulus alters the cycle decomposition of repeated doubling on \(\mathbb{Z}/m\mathbb{Z}\). The value \(m = 37\) produces the cleanest long orbit: length 36 from 1, with 2 acting as a primitive root modulo 37. Composites built from 37---notably \(111 = 3 \times 37\) and \(333 = 9 \times 37\)---inherit this long cycle length but fragment it according to the Chinese Remainder Theorem. Algebraic ``cleanness'' is therefore strongest at the repunit prime \textbf{37}, the prime factor of the length-3 repunit \(R_3 = 111\).

\textbf{Orderly progression of labels along the geometric orbit.} When the rotational step itself is \emph{coupled} to the modulus (\(\Delta\theta = m/\pi\)), labels advance with greater sequential coherence for certain composites. The value \textbf{111} frequently yields the highest scores for label progression and overall symmetry under this coupled regime. Here the modulus influences the \emph{flow of labels relative to the accumulating points}, rather than the intrinsic algebraic cycles alone. Composite structure and winding rate interact; the prime-mover number 111 is not merely decorative in this layer.

\textbf{Positional angle--label locking.} Only one labeling method creates a genuine \emph{statistical} dependence between label and actual geometric angle: explicit angular binning (\texttt{angle\_\allowbreak{}bin}). Under the default sequential method (\texttt{step\_\allowbreak{}index})---labels from the step count and its digital root or modular residue---normalized mutual information between labels and angles remains at or below chance once a proper null model (label shuffling with the same marginals) is applied.

Apparent visual symmetry in sequential plots therefore reflects \textbf{orderly sequential progression}, not a hidden locking of labels to preferred angles on the circle. That outcome is precisely what one should expect from an irrational rotation labeled by step index: the orbit has no preferred angles unless the labeling method itself introduces angular bins by construction.

\begin{figure}[htbp]
  \centering
  \includegraphics[width=0.92\textwidth,height=0.42\textheight,keepaspectratio]{fig9_5_modulus_invariants.png}
  \caption{Figure 9.5 --- Three invariants of modulus on an irrational rotation.}
  \label{fig:fig9-5-modulus-invariants}
\end{figure}

\begin{quote}\small\textit{Figure 9.5.* \textbf{(a)} Algebraic cleanness: \(m=37\), fixed step \(9/\pi\), sequential labels---long clean \(\times 2\) structure. \textbf{(b)} Orderly label flow: \(m=111\), coupled step \(m/\pi\), sequential labels---higher label-progression along the orbit. \textbf{(c)--(d)} Positive control pair at \(m=37\): sequential labeling yields excess NMI at or below the shuffle null; angular binning yields strongly positive excess NMI. Visual order under step-index is progression, not positional lock. *(Supporting generation: \texttt{vortex\_\allowbreak{}math} \(\cdot\) \texttt{assets/\allowbreak{}book\_\allowbreak{}figure\_\allowbreak{}three\_\allowbreak{}layers.\allowbreak{}png}.)}\end{quote}

\subsection{Controls and claim type}
\label{ch:ch09_algebras:controls-and-claim-type}

These layers were separated by controlled comparisons of mapping methods, step modes, and quantitative metrics: excess normalized mutual information (\textbf{exNMI}) against permutation baselines, label-progression scores, and angular uniformity. A positive control seals the analysis: when labels \emph{are} angle bins, exNMI is strongly positive and many standard deviations above the null; when labels are step indices, exNMI is not.

\textbf{Claim type.} The three-layer separation and control are a \textbf{Model} for how arithmetic moduli interact with geometric orbits in this book's ecosystem---experimentally reproducible, not a classical theorem of Hatcher \emph{Topology of Numbers}. Implementation and full tables: \textbf{Software fact} (\texttt{vortex\_\allowbreak{}math} repo; note in \texttt{notes/\allowbreak{}RESEARCH\_\allowbreak{}NOTE\_\allowbreak{}moduli.\allowbreak{}md}).

\subsection{Takeaway for the QGA spine}
\label{ch:ch09_algebras:takeaway-for-the-qga-spine}

The practical consequence is concise:

> \textbf{Modulus changes the algebraic and sequential-label invariants; it does not create angle--label locking on an irrational rotation unless you label by angle itself.}

This distinguishes synchronicity that arises from composite number structure---notably \textbf{111} as prime-mover composite \(3 \times 37\)---from stronger geometric claims that would require explicit angular construction. Algebra crowns 37; coupled progression frequently crowns 111 under the \(m/\pi\) regime; true positional locking belongs only to labeling methods that bin explicitly by angle.

For the rest of the book: when Chapters 5--8 speak of ``labels,'' ``classes,'' and ``Magic Islands,'' the same discipline applies---\textbf{name which invariant you are measuring}. Ideal class groups (?9.3--9.4) are algebraic invariants; flux scores and topograph diagrams are geometric or Model-level; statistical lock to a continuous base space is a third claim, and must be tested against a null, not inferred from visual order alone.

\begin{lstlisting}[style=qga]
# External supporting stack (not vendored in lib/)
~/Projects/vortex_math
  python src/main.py --family-37 --num-steps 200
  python src/main.py --resonance-scan --method step_index --num-steps 600
  python src/main.py --resonance-scan --method angle_bin --num-steps 600
\end{lstlisting}

\bigskip
\noindent\rule{\textwidth}{0.4pt}
\bigskip

\section{9.6 First computational labs}
\label{ch:ch09_algebras:9-6-first-computational-labs}

Helpers: \texttt{lib/\allowbreak{}quaternion\_\allowbreak{}algebra.\allowbreak{}py} \(\cdot\) \textbf{Appendix C ?C.4}.

\begin{itemize}
\item \textbf{9.A} \texttt{QuaternionAlgebra(-1,-1).\allowbreak{}ramified\_\allowbreak{}places()}.
\item \textbf{9.B} Hurwitz: 24 units, Euclidean, maximal.
\item \textbf{9.C} \texttt{left\_\allowbreak{}ideal\_\allowbreak{}class\_\allowbreak{}group} \(\rightarrow\) class number 1 (cited classical).
\item \textbf{9.D} Compare algebraic order 1 vs Model \texttt{class\_\allowbreak{}group\_\allowbreak{}analogue} order.
\item \textbf{9.E} Hilbert symbols table.
\item \textbf{9.F (optional, external)} Run the \texttt{vortex\_\allowbreak{}math} resonance control: \texttt{step\_\allowbreak{}index} vs \texttt{angle\_\allowbreak{}bin} and record exNMI signs.
\end{itemize}

\begin{lstlisting}[style=qga,language=python]
from lib.quaternion_algebra import QuaternionAlgebra, HurwitzOrder, left_ideal_class_group
print(QuaternionAlgebra(-1, -1).ramified_places())
print(HurwitzOrder().n_units(), left_ideal_class_group().order)
\end{lstlisting}

\bigskip
\noindent\rule{\textwidth}{0.4pt}
\bigskip

\section{Exercises}
\label{ch:ch09_algebras:exercises}

\textbf{9.A (hand).} Define a quaternion algebra over \(\mathbb{Q}\) and state what ramification means.

\textbf{9.B (hand).} Why is the Hurwitz order preferred over the Lipschitz order for most arithmetic work?

\textbf{9.C (code).} Complete Labs 9.A--9.B. Confirm 24 Hurwitz units and Euclidean / maximal flags.

\textbf{9.D (code).} Run Lab 9.C. What class number do you obtain for Hurwitz, and what does the \texttt{method} string say?

\textbf{9.E (code).} Complete Lab 9.D. Compare algebraic class number \(1\) with the Model \texttt{class\_\allowbreak{}group\_\allowbreak{}analogue} order. Why can they differ?

\textbf{9.F (Hatcher bridge).} In Hatcher Chapter 8, forms correspond to ideals. Sketch how ?9.4's dictionary might be made precise for a quaternion order (even if only conjecturally).

\textbf{9.G (Open Problem 6).} Using Hurwitz class number 1, what would a \emph{successful} OP6 composition law look like when restricted to configurations ``coming from'' principal ideals? What obstructions remain for non-principal classes in other algebras?

\textbf{9.H (forward).} Why will Chapter 10's observational validation benefit from a rigorous ideal-theoretic home for class-group analogues?

\textbf{9.I (software honesty).} Distinguish: (i) classical ideal class groups (\textbf{Theorem}), (ii) \texttt{left\_\allowbreak{}ideal\_\allowbreak{}class\_\allowbreak{}group} toy report (\textbf{software} citing theorem), (iii) Ch. 8 \texttt{class\_\allowbreak{}group\_\allowbreak{}analogue} (\textbf{Model}).

\textbf{9.J (modulus invariants).} In section 9.5, name the three layers (algebra / progression / angle lock). Which modulus wins algebraic cleanness? Which often wins progression under \(m/\pi\)? Why is positive exNMI under \texttt{angle\_\allowbreak{}bin} a \emph{control}, not a discovery about \(m=37\)?

\textbf{9.K (code, external).} If \texttt{\textasciitilde{}\allowbreak{}/\allowbreak{}Projects/\allowbreak{}vortex\_\allowbreak{}math} is available, run Lab 9.F and report exNMI under \texttt{step\_\allowbreak{}index} vs \texttt{angle\_\allowbreak{}bin} for \(m\in\{9,37,111\}\).

\bigskip
\noindent\rule{\textwidth}{0.4pt}
\bigskip

\section{Code and asset pointers}
\label{ch:ch09_algebras:code-and-asset-pointers}

\begin{lstlisting}[style=qga]
qga/lib/quaternion_algebra.py
  QuaternionAlgebra, hilbert_symbol,
  HurwitzOrder, LipschitzOrder,
  left_ideal_class_group, form_ideal_dictionary_entry

qga/lib/composition.py
qga/lib/hopf_lattice.py   # HURWITZ_UNITS shared with Ch. 3

# Supporting stack (external)
~/Projects/vortex_math/          # active development
qga/lib/vortex_math/README.md    # pointer into this book's tree
qga/notes/RESEARCH_NOTE_moduli.md
\end{lstlisting}

\textbf{Figures:} \texttt{scripts/\allowbreak{}generate\_\allowbreak{}ch9\_\allowbreak{}figures.\allowbreak{}py} \(\cdot\) Fig. 9.5 from \texttt{vortex\_\allowbreak{}math} book figure \textbf{Open problems:} OP6 (composition via ideal theory); OP1--OP3 for the geometric side of the dictionary. \textbf{Research note:} \texttt{notes/\allowbreak{}RESEARCH\_\allowbreak{}NOTE\_\allowbreak{}moduli.\allowbreak{}md}

\bigskip
\noindent\rule{\textwidth}{0.4pt}
\bigskip

\section{Looking ahead}
\label{ch:ch09_algebras:looking-ahead}

We now have the algebraic backbone---quaternion algebras, orders, and ideal class groups---that can host rigorous versions of the Model constructions in Chapters 5--8, together with a controlled lesson on \textbf{which invariant a modulus is measuring} when discrete labels meet continuous orbits. In \textbf{Chapter 10} we return to the observational layer: statistical protocols for \(350/\pi\), the \(Z\mapsto\) map, Magic Islands, and any class-group predictions that become precise enough to test---always distinguishing algebraic claims, Model dictionaries, and empirical lock statistics.

With the ideal-theoretic foundation in place, the full construction is ready for observational scrutiny.

\bigskip
\noindent\rule{\textwidth}{0.4pt}
\bigskip


